\chapter{逐篇完整性与来源审查}
\label{app:audit}
\section{审查标准与本版交付边界}
本版的“解释完整”分为两个层面：\textbf{教学闭合}要求读者知道输入、计算对象、更新、结束、算例和结论条件；\textbf{原文可复现}还要求原始模型、全部超参数、数据协议和实现细节已获得并核对。前者可以通过明确标记的教学变体补足，后者不能由自行补写冒充。任何本地数值验算都不等于复现原论文实验。

\begin{longtable}{p{0.08\textwidth}p{0.24\textwidth}p{0.34\textwidth}p{0.24\textwidth}}
\toprule 编号 & 当前来源范围 & 本版重点补齐的算法内容 & 仍需保留的边界\\\midrule\endhead
S1 & 公开预印本全文 & 数据路由、相量回归及协方差算例 & 教程，不是单一算法\\
S2 & 出版元数据；先前预览记录 & 三阶段公共模块的独立推导 & 未得全文；不可称原算法复现\\
S3 & 作者机构全文 & 投影设计行、QP、分组状态、回溯、评分、完整树算例 & 原文若干索引/符号约定须核对；有限候选\\
S4 & 出版摘要 & 高斯潜树EM的E/M更新、手算、BIC及停止 & 教学模型非已核实原模型；未得全文\\
S5 & 本轮补得出版全文 & 多源EM、折扣、证据融合、决策与数值轨迹 & 缺失权重和置信表达的文本/公式差异\\
S6 & 出版记录 & 可确认任务、应核对的训练/评估接口 & 未得全文；不虚构网络结构\\
S7 & 公开预印本全文 & 量化设计、受约束估计、支持恢复条件与算例 & 教学求解器单独标记；理论假设接口\\
S8 & 作者机构全文 & 参考序列、相似度、适应阈值、归属判断 & 现场覆盖与准确率含义不同\\
S9 & 作者机构全文 & WLS、残差投影、标准化及诊断路径 & 异常残差不能唯一定位错误原因\\
S10 & 作者公开全文 & 四线模型、灵敏度回归、包络分配步骤 & 预测/模型误差与运行限制需分开\\
S11 & 公开预印本全文 & 椭球、矩形、去冗余、扩张及二维算例 & 线性模型鲁棒包含不等于AC全域安全\\
S12 & 公开预印本全文 & 探测矩阵、等值集合、递归恢复与树算例 & 全节点/部分量测决定不同恢复对象\\
S13 & 公开预印本全文 & 投影LS、二值化、MILP和连通约束 & 有界变量及候选线路先验\\
P1 & 公开预印本全文 & Laplacian目标、ADMM更新、投影与停止 & 等价记号整理；树后处理与凸求解分开\\
P2 & 出版全文 & 特征提取、祖先记录、传递约简 & 上游可见性依赖电路与检测条件\\
P3 & 出版全文 & 检测、时延、反射、长度换算 & 高频模型及传播参数；非静态R/X算法\\
P4 & 出版全文 & 记录包含、闭包、缺失与方向未决处理 & 逻辑蕴含和经验补全分别标记\\
P5 & 公开预印本全文 & 相关解码、Hankel/ERA及端口模型 & 端口等效实现不唯一对应物理树\\
S14 & 公开全文 & KL变分、归一化、温度和有限候选算例 & 广义后验权重非自动频率覆盖\\
S15 & 公开预印本全文 & 训练密度、profile分母、比率及保守界 & 独立划分、模型正确性、优化界方向\\
S16 & 作者公开全文 & 树空间状态、拓扑提案、边长MH与数值轮次 & 几何/接受式的版本级审查\\
S17 & 公开预印本全文 & 校准分数、有限样本分位数和集合构造 & 交换性、样本单位、预测/拓扑目标\\
\bottomrule
\end{longtable}

\section{原稿中需要改进的解释，现在如何闭合}
旧稿关于“当前模型能否迁移”的讨论较多，而算法描述有时停在目标层。本版新增的关键连接包括：S3每个量测如何写入设计矩阵；RG合并后活动集、距离及图如何同步变化；S12如何从等值集合逐层生成子树；P1如何从增广目标得到各个更新块；P5如何从相关序列构造ERA矩阵；S4的E步为何要保留二阶矩；S15的优化界为何必须朝特定方向；S16怎样提出并接受树与边长变化；S17怎样处理超出有限样本最大索引的分位数。各章给出具体操作及数值例，而非只列算法名称。

审查后仍不应把全书称为“22篇全部完整复现”。S2、S4、S6的缺口属于来源限制；其余文献本版也未重跑作者实验或核对全部作者代码。读取正文、证明教学推导、检查数值例和编译排版，是本次已经完成或执行的验证层次；精确实验复现是另一项工作。

\section{公式冲突怎样处理：采用固定定义重新验证}
文献中出现不一致时，先检查是否因向量化、符号方向或单位不同造成，再对照原文页面，最后用维数、特殊例子或目标的一阶条件验证。本版不把文本抽取异常直接宣布为原文错误；只有核对页面后才记录版本级问题。下表集中列出对理解有实质影响的接口，各专题节仍给出具体推导。

\begin{longtable}{p{0.11\textwidth}p{0.36\textwidth}p{0.43\textwidth}}
\toprule 来源 & 所读版本中需要检查的表达 & 本版采用的可验证处理\\\midrule\endhead
S1 & 完整关联矩阵分块后的根列公式 & 从 $[a_0\ A]\mathbf1=0$ 直接得 $a_0=-A\mathbf1$，不依赖错置的逆矩阵。\\
S3 & 原文第3.2节父子方向文字与 $\Phi_{ijk}=D_{ik}-D_{jk}$ 的定义 & 若子 $i$ 经父 $j$ 到外部 $k$，必有 $\Phi=+D_{ij}$。按路径等式和根方向判断，记录与文字的符号差异。\\
S3 & 式(24)、(26)的平均集合/分母，以及相对分数零值 & 明确外部集合 $A\setminus\{i,j\}$ 和分母；零值/负值规则不冒充原文已规定。\\
S5 & 缺失权重及置信度的文字与打印式不完全一致 & 分别解释实际式子与文字目标，教学约定写在算法旁边。\\
P1 & 局部变量副本与乘子编号存在互换 & 以一致增广拉格朗日重新导出更新，标为等价代数整理。\\
S13 & 矩阵逆与维数，以及树约束的适用条件 & 从最小二乘消元重推；连接数和非孤立不足以替代连通性，流约束须明确方向。\\
S16 & 树空间乘积距离与CAT(0)解释；MH式出现 $\max$ & 几何结论按所用乘积度量检查；MH概率由详细平衡推导为 $\min\{1,r\}$。\\
S11 & 上下界符号及去冗余的删除顺序 & 固定 $f=(u,-l)$ 的明确记号；采用逐行删除复检保持多面体不变。\\
\bottomrule
\end{longtable}
这些记录不意味着相应方法整体失效。特别是启发式评分是否实用，需要其适用数据上的实验；一个公式或符号接口的核验结果，不应被扩大成对整篇论文结论的否定。

\section{来源、版本和复现文件}
每篇文献的完整书目信息在文末，原始来源的URL、版本及正文定位在随源文件附带的 \path{sources_core.md}、\path{sources_engineering.md}、\path{sources_probing.md} 和 \path{sources_statistical.md} 中。逐专题审查记录位于 \path{audit_*.md}。正文以章节和算法编号组织；来源文件保留原论文的节号、式号，避免本书重新编号后无法追溯。

主文件为 \path{main.tex}，章节在 \path{chapters/}，书目在四个 \path{refs_*.bib}。\path{build.ps1} 使用 XeLaTeX 与 BibTeX 自动构建，\path{assemble_standalone.py} 生成嵌入书目的单文件版本。数值核验脚本只检查教学例的代数一致性；不作为原论文实验成绩。旧版目录 \path{docs/literature_detailed_20260910/} 保留，便于对照研究审查稿与本教材稿。
